\begin{circuitikz}[x=1cm,y=1cm,>=stealth]
	
	% Grundschrift
	\tikzstyle{every node}=[font=\fontsize{13pt}{16pt}\selectfont]
	
	% Kleine Schrift fuer Pfeilbeschriftungen
	\tikzset{
		pfeiltext/.style={
			font=\fontsize{6.5pt}{8pt}\selectfont,
			fill=white,
			fill opacity=1,
			text opacity=1,
			inner xsep=0.025cm,
			inner ysep=0.015cm,
			align=center
		}
	}
	
	% Gesamtsystem
	\draw [line width=2pt, dashed] (1.0,16.75) rectangle (16.5,-1.125);
	
	\node [
	font=\fontsize{25.0pt}{32.5pt}\selectfont,
	fill=white,
	inner xsep=0.080cm,
	inner ysep=0.085cm,
	rounded corners=0.020cm
	] at (9.125,16) {\textbf{Gesamtsystem}};
	
	% Maschine
	\draw [
	color={rgb,255:red,246; green,27; blue,27},
	draw opacity=1,
	line width=2pt,
	dashed
	] (4.375,6.625) rectangle (14.125,-0.125);
	
	\node [
	font=\fontsize{18.2pt}{23.7pt}\selectfont,
	color={rgb,255:red,253; green,17; blue,17},
	fill=white,
	text opacity=1,
	inner xsep=0.080cm,
	inner ysep=0.085cm,
	rounded corners=0.020cm
	] at (9.125,5.95) {\textbf{Maschine}};
	
	\draw [line width=1.2pt] (5,5.625) rectangle (8,3.5);
	\draw [line width=1.2pt] (10.375,5.625) rectangle (13.25,3.5);
	
	\node [
	font=\fontsize{13.1pt}{17.0pt}\selectfont,
	inner xsep=0.080cm,
	inner ysep=0.085cm,
	rounded corners=0.020cm
	] at (6.5,4.625) {\textbf{Heizdraht}};
	
	\node [
	font=\fontsize{13.1pt}{17.0pt}\selectfont,
	inner xsep=0.080cm,
	inner ysep=0.085cm,
	rounded corners=0.020cm
	] at (11.75,4.5) {\textbf{Motoren}};
	
	% Steuerung
	\draw [
	color={rgb,255:red,14; green,170; blue,248},
	draw opacity=1,
	line width=2pt,
	dashed
	] (4.375,15.25) rectangle (13.875,7.5);
	
	\node [
	font=\fontsize{18.2pt}{23.7pt}\selectfont,
	color={rgb,255:red,95; green,170; blue,242},
	fill=white,
	text opacity=1,
	inner xsep=0.080cm,
	inner ysep=0.085cm,
	rounded corners=0.020cm
	] at (9.125,14.75) {\textbf{Steuerung}};
	
	\draw [line width=1.2pt] (6.75,14.125) rectangle (11.375,13);
	
	\node [
	font=\fontsize{18.2pt}{23.7pt}\selectfont,
	inner xsep=0.080cm,
	inner ysep=0.085cm,
	rounded corners=0.020cm
	] at (9.125,13.5) {\textbf{Computer}};
	
	\draw [line width=1.2pt] (5.125,11.625) rectangle (8,10.25);
	\draw [line width=1.2pt] (10.25,11.625) rectangle (13.125,10.25);
	\draw [line width=1.2pt] (5.125,9.375) rectangle (8,8.125);
	\draw [line width=1.2pt] (10.25,9.375) rectangle (13.125,8.125);
	
	\node [
	font=\fontsize{9.2pt}{11.5pt}\selectfont,
	inner xsep=0.040cm,
	inner ysep=0.040cm,
	rounded corners=0.020cm
	] at (6.5,10.875) {\textbf{Mikrocontroller}};
	
	\node [
	font=\fontsize{9.2pt}{11.5pt}\selectfont,
	inner xsep=0.040cm,
	inner ysep=0.040cm,
	rounded corners=0.020cm
	] at (11.625,10.75) {\textbf{CNC-Shield}};
	
	\node [
	font=\fontsize{10.8pt}{14.1pt}\selectfont,
	inner xsep=0.080cm,
	inner ysep=0.085cm,
	rounded corners=0.020cm
	] at (6.5,8.75) {\textbf{MOSFET}};
	
	\node [
	font=\fontsize{9.2pt}{11.5pt}\selectfont,
	inner xsep=0.040cm,
	inner ysep=0.040cm,
	rounded corners=0.020cm
	] at (11.625,8.625) {\textbf{Motor-Treiber}};
	
	% Sensoren und Bedienelemente
	\draw [line width=1.2pt] (10.375,2.375) rectangle (13.375,0.5);
	
	\node [
	font=\fontsize{13.1pt}{17.0pt}\selectfont,
	inner xsep=0.080cm,
	inner ysep=0.085cm,
	rounded corners=0.020cm
	] at (11.875,1.5) {\textbf{Endschalter}};
	
	\draw [line width=1.2pt] (5,2.375) rectangle (8,0.5);
	
	\node [
	font=\fontsize{11.9pt}{15.5pt}\selectfont,
	inner xsep=0.080cm,
	inner ysep=0.085cm,
	rounded corners=0.020cm
	] at (6.5,1.5) {\textbf{Druckknöpfe}};
	
	% Serielle Kommunikation
	\draw [line width=0.9pt, <->] (7.5,13) -- (7.5,11.625)
	node[
	pos=0.5,
	left,
	pfeiltext
	]{\textbf{serielle}\\\textbf{Kommunikation}};
	
	% Shield Header
	\draw [line width=0.9pt, ->] (8,11.125) -- (10.25,11.125)
	node[
	pos=0.5,
	above,
	pfeiltext
	]{\textbf{Shield-Header}};
	
	% STEP DIR ENABLE
	\draw [line width=0.9pt, ->] (11.625,10.25) -- (11.625,9.375)
	node[
	pos=0.5,
	right,
	pfeiltext
	]{\textbf{STEP, DIR}\\\textbf{ENABLE}};
	
	% Motorstrom
	\draw [line width=0.9pt, ->] (11.75,8.125) -- (11.75,5.625)
	node[
	pos=0.35,
	right,
	pfeiltext
	]{\textbf{Motorstrom}};
	
	% Laststrom
	\draw [line width=0.9pt, ->] (6.375,8.125) -- (6.375,5.625)
	node[
	pos=0.35,
	left,
	pfeiltext
	]{\textbf{Laststrom}};
	
	% PWM-Signal
	\draw [line width=0.9pt, ->] (6.5,10.25) -- (6.5,9.375)
	node[
	pos=0.5,
	right,
	pfeiltext
	]{\textbf{PWM-Signal}};
	
	% Endschalter-Rueckmeldung
	\draw [line width=0.9pt] (11.875,0.5) -- (11.875,0.05);
	\draw [line width=0.9pt] (11.875,0.05) -- (2.15,0.05);
	\draw [line width=0.9pt] (2.15,0.05) -- (2.15,10.95);
	\draw [line width=0.9pt, ->] (2.15,10.95) -- (5.125,10.95);
	
	\node[
	pfeiltext
	] at (3.25,11.35) {\textbf{Endschalter-}\\\textbf{Rückmeldung}};
	
	% Druckknopf-Signal
	\draw [line width=0.9pt] (5,1.5) -- (2.75,1.5);
	\draw [line width=0.9pt] (2.75,1.5) -- (2.75,10.45);
	\draw [line width=0.9pt, ->] (2.75,10.45) -- (5.125,10.45);
	
\node[
pfeiltext
] at (3.35,1.05) {\textbf{digitales}\\\textbf{Eingangssignal}};
	
\end{circuitikz}